% \iffalse meta-comment
%
% Copyright (C) 2026 by Laurent Coste
%
% This file may be distributed and/or modified under the
% conditions of the LaTeX Project Public License, either
% version 1.3 of this license or (at your option) any later
% version. The latest version of this license is in:
%
%    http://www.latex-project.org/lppl.txt
%
% and version 1.3 or later is part of all distributions of
% LaTeX version 2005/12/01 or later.
%
% \fi
%
% \iffalse
%<*driver>
\ProvidesFile{seqenseigne-style6e.dtx}
%</driver>
%<package>\NeedsTeXFormat{LaTeX2e}
%<package>\ProvidesPackage{seqenseigne-style6e}
%<*package>
    [2026/10/20 v0.8 seqenseigne : style des documents de 6e]
%</package>
%
%<*driver>
\documentclass{ltxdoc}
\usepackage[utf8]{inputenc}
\usepackage[T1]{fontenc}
\IfFileExists{french.ldf}{\usepackage[french]{babel}}{}
\usepackage[margin=2.2cm,left=4cm]{geometry}
\usepackage{seqenseigne-style6e}
\EnableCrossrefs
\CodelineIndex
\RecordChanges
\begin{document}
  \DocInput{seqenseigne-style6e.dtx}
  \PrintChanges
  \PrintIndex
\end{document}
%</driver>
% \fi
%
% \CheckSum{0}
%
% \changes{v0.8}{2026/10/20}{Astuce dans un énoncé : \cs{seqAstuceSix} (ampoule,
%   « Astuce » en gras, texte en petits caractères, sans cadre ; élève et corrigé)}
% \changes{v0.7}{2026/10/20}{Marqueurs de l'exercice : \cs{seqMaterielSix},
%   \cs{seqOrdinateurSix} et \cs{seqCalculatriceSix} (pictogramme et nom de
%   l'outil, calculatrice autorisée ou interdite), dans l'en-tête de
%   \texttt{seqExoSix} et de \texttt{seqActiviteSix}}
% \changes{v0.5}{2026/10/19}{Niveau paramétrable : \cs{seqNiveauCourt} et
%   \cs{seqNiveauExposant} (6e par défaut) remplacent « 6e » écrit en dur dans
%   la couverture du livret, l'identification d'évaluation, les en-têtes de
%   travail personnel et d'activité complémentaire}
% \changes{v0.1}{2026/10/18}{Prototype : livret, calcul mental, évaluation, fiche de synthèse}
% \changes{v0.2}{2026/10/18}{Activité d'introduction, fiche de résumé à compléter,
%   fiche d'activité complémentaire, planches à découper au massicot}
% \changes{v0.4}{2026/10/18}{Intégration des macros \cs{seqLoc} des lots S01 à S03}
% \changes{v0.3.1}{2026/10/18}{\cs{seqACompleter} utilisable en mode mathématique}
% \changes{v0.3}{2026/10/18}{Mise en \texttt{.dtx} ; réponses unifiées (version élève
%   ou corrigée par une seule bascule) ; abandon du cadre « À retenir » ;
%   réglages \textsf{tkz-euclide} ; en-tête de travail personnel avec délai}
%
% \GetFileInfo{seqenseigne-style6e.dtx}
%
% \title{Le module \textsf{seqenseigne-style6e}\thanks{Ce document correspond à
%   \textsf{seqenseigne-style6e}~\fileversion, du \filedate.}}
% \author{Laurent Coste}
% \maketitle
%
% \begin{abstract}
% Ce module du paquet \textsf{seqenseigne} donne aux documents de 6\textsuperscript{e}
% le visuel des documents de l'équipe de mathématiques du collège : livrets d'exercices,
% fiches de mise en route, fiches de cours, évaluations, travaux personnels, fiches
% d'activité complémentaire. Il ne modifie aucune macro des autres modules et
% n'est chargé que pour le cycle~3.
% \end{abstract}
%
% \tableofcontents
%
% \section{Principes}
%
% \paragraph{Une seule source, deux versions.} Toutes les zones de réponse passent
% par les macros de réponse (section~\ref{sec:reponses}). Elles reçoivent la
% réponse attendue en argument. La version élève affiche des pointillés ou une case
% vide. La version corrigée, ou complète pour une fiche de cours, affiche la réponse
% en rouge gras. On bascule de l'une à l'autre avec \cs{seqVersionCorrigee} ou
% \cs{seqVersionEleve}, sans réécrire le document.
%
% \paragraph{Correspondance avec l'appli.} Les statuts des livrets de l'équipe
% deviennent les séries de l'appli : (O) $\to$ série fondamentale, (A) $\to$ série
% avancée, [R] $\to$ remédiation de l'exercice fondamental. Le bandeau de capacité
% est le titre de la partie ou de l'objectif, émis par l'appli.
%
% \paragraph{Chargement.} Le module est chargé par l'option \texttt{sixieme} du
% paquet, après les modules de base : |\usepackage[sixieme]{seqenseigne}|. Il peut
% aussi être chargé seul, pour des essais.
%
% \paragraph{Conventions du paquet.} Les macros publiques sont préfixées par
% \cs{seq}, les macros internes par \cs{seq@}. On n'utilise jamais
% \cs{robustify}\verb|{\par}|. Les macros déjà définies ailleurs ne sont pas
% redéfinies.
%
% \section{Utilisation}
%
% \subsection{Niveau}
% \DescribeMacro{\seqNiveauCourt}\DescribeMacro{\seqNiveauExposant}
% Le niveau des élèves s'affiche sur la couverture du livret et dans les champs
% d'identification. \cs{seqNiveauCourt} vaut \texttt{6e} et \cs{seqNiveauExposant}
% vaut 6\textsuperscript{e} par défaut. L'appli les redéfinit avec le nom court du
% niveau du référentiel, par exemple |\renewcommand{\seqNiveauCourt}{5e}| et
% |\renewcommand{\seqNiveauExposant}{5\textsuperscript{e}}| : le modèle de
% documents ne dépend pas du niveau.
%
% \subsection{Couleurs}
% Le module définit les couleurs \texttt{seqSixMauve}, \texttt{seqSixMauveBord},
% \texttt{seqSixRouge}, \texttt{seqSixBleu}, \texttt{seqSixVert},
% \texttt{seqSixPointilles}, \texttt{seqSixGrisTexte} et \texttt{seqSixEnteteTab}.
% Il les relève sur les documents 2026 de l'équipe.
%
% \subsection{Réponses}\label{sec:reponses}
% \DescribeMacro{\seqVersionEleve}\DescribeMacro{\seqVersionCorrigee}
% Ces deux macros basculent tout ce qui suit entre la version élève (par défaut) et la
% version corrigée ou complète. \cs{seqVersionComplete} et \cs{seqVersionACompleter} sont
% des synonymes, réservés aux fiches de cours.
%
% \DescribeMacro{\seqACompleter}
% \cs{seqACompleter}\oarg{largeur}\marg{réponse} affiche des pointillés en ligne dans la
% version élève, et la réponse dans la version corrigée. Sans \meta{largeur}, les
% pointillés prennent la largeur de la réponse plus 1,2\,em : la mise en page reste
% identique d'une version à l'autre.
%
% \DescribeMacro{\seqReponse}
% \cs{seqReponse}\oarg{n}\marg{réponse} réserve $n$ lignes de pointillés pleine largeur
% (une par défaut). La version corrigée écrit la réponse sur ces lignes.
%
% \DescribeMacro{\seqCase}
% \cs{seqCase}\oarg{largeur}\oarg{réponse} trace une case à remplir dans le texte.
%
% \DescribeMacro{\seqCaseCocher}
% \cs{seqCaseCocher}\oarg{x}\marg{texte} trace une case à cocher. Avec l'option
% \texttt{x}, la case est cochée dans la version corrigée.
%
% \DescribeMacro{\seqTrous}\DescribeMacro{\seqPointilles}
% \cs{seqTrous}\oarg{longueur} et \cs{seqPointilles}\oarg{n}\oarg{couleur} sont les briques
% de base, sans réponse : on les utilise pour un champ d'identification (nom, date).
%
% \subsection{Livret d'exercices}
% \DescribeMacro{\seqCouvertureLivretSix}
% \cs{seqCouvertureLivretSix}\marg{titre}\marg{liste} produit la page de garde : bloc
% NOM, Prénom et classe, titre et niveau, puis la liste des objectifs du chapitre.
% La \meta{liste} est de la forme \verb|C1/{texte},C2/{texte}|.
%
% \DescribeMacro{\seqPartieCouvertureSix}
% Depuis la v0.6, la forme étoilée \cs{seqCouvertureLivretSix*}\marg{titre}\marg{parties}
% range les objectifs par partie : \meta{parties} est une suite de
% \cs{seqPartieCouvertureSix}\marg{titre de la partie}\marg{liste}, avec la même
% \meta{liste} que ci-dessus. Chaque partie a son titre, puis ses objectifs avec un
% trait vertical à gauche, de la couleur \texttt{seqSixTheme} (couleur du thème de la
% séquence, |\definecolor{seqSixTheme}{HTML}{7326A0}| par exemple ; bleu du module
% par défaut).
%
% \DescribeMacro{\seqBandeauCapacite}
% \cs{seqBandeauCapacite}\marg{code}\marg{intitulé} ouvre la section d'un objectif.
%
% \DescribeEnv{seqActiviteSix}
% L'environnement |seqActiviteSix|\marg{titre} met en forme l'activité
% d'introduction, au début de la capacité. Elle porte une étiquette bleue numérotée et
% un filet à gauche, sans carré de numéro : ce carré reste réservé aux exercices.
%
% \DescribeEnv{seqFicheResumeSix}
% L'environnement |seqFicheResumeSix|\marg{titre} contient la fiche de résumé de
% l'objectif. Dans le livret, en version élève, il ne reste que quelques mots à
% compléter. La même fiche, en version complète, sert de fiche de cours.
%
% \DescribeEnv{seqExoSix}
% L'environnement |seqExoSix|\oarg{statut}\oarg{sous-titre} contient un exercice
% numéroté. Le numéro est inscrit dans un carré mauve, suivi du statut : \texttt{O},
% \texttt{A} ou \texttt{R}. Le statut \texttt{O} est la valeur par défaut. La
% numérotation est continue sur tout le livret.
%
% \subsection{Marqueurs de l'exercice (v0.7)}
% \DescribeMacro{\seqMaterielSix}
% \cs{seqMaterielSix}\marg{outil}\marg{calculatrice} se place dans le titre d'un
% exercice ou d'une activité, avant le code. Elle affiche ce que l'élève doit
% savoir avant de commencer : l'ordinateur et l'outil, puis la calculatrice.
% Un argument vide n'affiche rien. La mise en forme est sobre, en petits
% caractères gris foncé, et reste lisible en noir et blanc.
%
% \DescribeMacro{\seqOrdinateurSix}
% \cs{seqOrdinateurSix}\marg{outil} affiche un écran suivi de « Ordinateur --- outil »,
% par exemple « Ordinateur --- Scratch ». Si l'outil est vide, seul « Ordinateur »
% est écrit.
%
% \DescribeMacro{\seqCalculatriceSix}
% \cs{seqCalculatriceSix}\marg{autorisee ou interdite} affiche une calculatrice suivie
% de « calculatrice autorisée », ou une calculatrice barrée suivie de « calculatrice
% interdite ».
%
% \DescribeMacro{\seqPictoOrdinateurSix}\DescribeMacro{\seqPictoCalculatriceSix}
% Les deux pictogrammes seuls ; la forme étoilée \cs{seqPictoCalculatriceSix*}
% barre la calculatrice.
%
% \subsection{Astuce (v0.8)}
% \DescribeMacro{\seqAstuceSix}
% \cs{seqAstuceSix}\marg{texte} ajoute une astuce à la fin d'un énoncé, par exemple
% « Pour mieux comprendre ce que fait le programme, tu peux le construire dans Scratch
% et le tester. » Elle s'affiche dans un paragraphe à part : une ampoule, le mot
% « Astuce » en gras, puis le texte en petits caractères, sans cadre. Les lignes
% suivantes s'alignent sur le mot « Astuce ». L'astuce est visible dans la version
% élève et dans la version corrigée. C'est une suggestion, pas une consigne :
% un exercice qui porte une astuce n'est pas marqué « ordinateur ».
%
% \DescribeMacro{\seqPictoAstuceSix}
% Le pictogramme seul : une ampoule au trait, lisible en noir et blanc.
%
% \subsection{Mises en route}
% \DescribeEnv{seqCalculMental}
% L'environnement |seqCalculMental|\marg{numéro}\marg{titre}\marg{barème}\oarg{colonnes}
% produit une fiche de séance : le bandeau « CALCUL MENTAL \meta{numéro} -- \meta{titre} »,
% les items rangés ligne par ligne, puis le bloc SCORE.
%
% \DescribeMacro{\seqItemCM}
% \cs{seqItemCM}\marg{calcul}\oarg{résultat} donne un item. Le résultat ne s'affiche que
% dans la version corrigée.
%
% \DescribeMacro{\seqScore}
% \cs{seqScore}\marg{n} affiche le bloc d'autoévaluation « SCORE : \dots\ / $n$ ».
%
% \subsection{Planches à découper au massicot}
% \DescribeEnv{seqPlancheFiches}\DescribeEnv{seqBande}
% |seqPlancheFiches|\marg{n} place $n$ fiches différentes par page A4. Chaque fiche est
% un environnement |seqBande|, et un trait de coupe sépare deux fiches.
%
% \DescribeEnv{seqCopies}
% |seqCopies|\marg{n} reproduit $n$ fois le même contenu sur une page, par exemple pour
% un cours ou un test en A5 avec $n=2$. Les compteurs du module repartent à la même
% valeur pour chaque copie.
%
% \subsection{Évaluations}
% \DescribeMacro{\seqIdentEleve}\DescribeMacro{\seqTitreEvalSix}
% \DescribeMacro{\seqBaremeCompetences}\DescribeMacro{\seqExoEval}\DescribeMacro{\seqItemRond}
% Ces macros reprennent le gabarit des évaluations 2026 de l'équipe :
% \begin{itemize}
% \item une ligne d'identification ;
% \item un titre et un sous-titre, avec la séquence, le sujet A ou B et la calculatrice ;
% \item le barème par compétence (\verb|\seqBaremeCompetences{C1/3,C2/5,Bonus/2}{20}|) ;
% \item les titres d'exercice « Exercice $n$ \meta{titre} (\meta{points} --- \meta{compétence}) » ;
% \item les items cerclés.
% \end{itemize}
%
% \subsection{Fiches de cours, travaux, activités complémentaires}
% \DescribeMacro{\seqTitreFicheSix}\DescribeMacro{\seqSectionFicheSix}
% \DescribeMacro{\seqMotCle}\DescribeEnv{seqEncadreDefSix}
% Fiche de cours : titre vert souligné, sections bleues soulignées, mots-clés rouges,
% encadré rouge pour les définitions.
%
% \DescribeMacro{\seqTitreTravailSix}
% \cs{seqTitreTravailSix}\oarg{délai}\marg{numéro}\marg{titre}\marg{contexte} produit
% l'en-tête d'un travail personnel. Ce travail est donné à toute la classe, à faire à
% la maison, et il est noté. Le \meta{délai} s'affiche à droite, par exemple
% « à rendre le \dots ».
%
% \DescribeEnv{seqEncadreIntroSix}\DescribeMacro{\seqPartieSix}
% Encadré bleu de mise en situation ; titres de parties « Partie A --- \dots ».
%
% \DescribeMacro{\seqTitreComplementSix}
% \cs{seqTitreComplementSix}\marg{titre}\marg{contexte} produit l'en-tête de la fiche
% d'activité complémentaire. Cette fiche est un document à part, distribué aux élèves
% qui ont terminé.
%
% \subsection{Figures}
% \DescribeMacro{\seqDemiDroite}
% \cs{seqDemiDroite}\oarg{longueur}\marg{origine}\marg{unité}\marg{nb}\marg{sous-graduations}\oarg{points}
% trace une demi-droite graduée. Les \meta{points} sont donnés sous la forme
% \verb|1.3/A,2.7/B|.
%
% \DescribeMacro{\seqFiguresSix}
% Pour les figures de géométrie, on utilise \textsf{tkz-euclide}, qui fait apparaître
% les propriétés : parallèles, perpendiculaires, codages de longueurs et d'angles,
% symétriques. \cs{seqFiguresSix} règle les points et les codages au style des
% documents de l'équipe : points en croix et lettres en gras. Le module l'applique
% au début du document quand \textsf{tkz-euclide} est chargé, c'est-à-dire avec
% l'option \texttt{geometrie} du paquet : |\usepackage[sixieme,geometrie]{seqenseigne}|.
%
% \subsection{Géométrie plane (v0.6)}
% Ces macros viennent de la conversion de la séquence S05 « Géométrie plane ».
% Elles gardent leur préfixe \cs{seqLoc} : les énoncés du lot les utilisent
% tels quels.
%
% \DescribeMacro{\seqLocCadreTrace}
% \cs{seqLocCadreTrace}\oarg{hauteur}\marg{figure} réserve une zone de tracé :
% un cadre gris clair de la largeur de la ligne et de la \meta{hauteur} donnée
% (5\,cm par défaut). En version élève, le cadre est vide : l'élève y construit
% la figure. En version corrigée, la \meta{figure} (en général un
% \texttt{tikzpicture}) est centrée dans le même cadre. La mise en page est
% identique dans les deux versions.
%
% \DescribeMacro{\seqLocAngle}
% \cs{seqLocAngle}\marg{ABC} écrit l'angle géométrique $\widehat{ABC}$, avec
% son chapeau, en mode texte comme en mode mathématique.
%
% \DescribeMacro{\seqLocRapporteur}
% \cs{seqLocRapporteur}\oarg{rayon}\marg{centre}\marg{orientation}, dans un
% \texttt{tikzpicture}, dessine un demi-rapporteur semi-transparent à double
% graduation (rayon de 3 unités par défaut). Son centre est au point
% \meta{centre}, par exemple |{(2,0)}| ; la graduation 0 de l'échelle
% extérieure est dans la direction \meta{orientation}, en degrés. L'échelle
% extérieure va de 0 à 180 dans le sens inverse des aiguilles d'une montre,
% l'échelle intérieure (en rouge) de 180 à 0. On voit l'angle à travers le
% rapporteur. Il remplace \cs{tkzProtractor}, peu lisible une fois imprimé.
%
% Le style TikZ \texttt{seqLocCorr} (rouge, trait de 0,8\,pt) est celui des
% tracés du corrigé. On l'emploie avec la bascule élève\,/\,corrigé :
% |\seqLocSiCorrige{\begin{scope}[seqLocCorr] ... \end{scope}}{}|.
%
% \subsection{Pré-algèbre (v0.6)}
% Ces macros viennent de la conversion de la séquence S04 « Pré-algèbre » et gardent
% leur préfixe \cs{seqLoc}.
%
% Schémas en barres, dans un \texttt{tikzpicture} (unité : le cm). Les styles TikZ
% \texttt{seqLocBarre} (barre bleue), \texttt{seqLocBarreR} (rouge),
% \texttt{seqLocBarreV} (verte) et \texttt{seqLocBarreQ} (pointillée : la quantité
% cherchée) donnent l'aspect des barres.
%
% \DescribeMacro{\seqLocBloc}
% \cs{seqLocBloc}\oarg{style}\marg{x}\marg{y}\marg{longueur}\marg{étiquette} trace un
% bloc de barre de 0,6\,cm de haut, de coin inférieur gauche $(x,y)$, avec
% l'\meta{étiquette} au centre (style \texttt{seqLocBarre} par défaut).
%
% \DescribeMacro{\seqLocAccolade}
% \cs{seqLocAccolade}\oarg{position}\marg{x1}\marg{x2}\marg{y}\marg{étiquette} trace une
% accolade horizontale de $x_1$ à $x_2$ à l'ordonnée $y$, sous la barre
% (\texttt{below}, par défaut) ou au-dessus (\texttt{above}), avec son étiquette.
%
% \DescribeMacro{\seqLocBalance}
% \cs{seqLocBalance}\oarg{demi-largeur}\marg{plateau gauche}\marg{plateau droit}
% dessine une balance de Roberval en équilibre (demi-largeur du fléau : 2,4\,cm par
% défaut). Les plateaux reçoivent du texte ou des objets posés côte à côte.
%
% \DescribeMacro{\seqLocSac}\DescribeMacro{\seqLocMasse}\DescribeMacro{\seqLocObjet}
% Objets à poser sur la balance : \cs{seqLocSac}\marg{nom} (sac ou boîte de masse
% inconnue), \cs{seqLocMasse}\marg{valeur} (masse marquée),
% \cs{seqLocObjet}\oarg{couleur}\marg{texte} (objet rond : fruit, bille\dots ; rouge par
% défaut).
%
% \StopEventually{}
%
% \section{Implémentation}
%
%    \begin{macrocode}
%<*package>
\RequirePackage{xcolor}
\RequirePackage{tcolorbox}
\tcbuselibrary{skins,breakable}
\RequirePackage{tikz}
\usetikzlibrary{calc}
\RequirePackage{enumitem}
\RequirePackage{xparse}
\RequirePackage{array}
\RequirePackage{pifont}
\RequirePackage{graphicx}
%    \end{macrocode}
%
% \subsection{Niveau}
% \begin{macro}{\seqNiveauCourt}\begin{macro}{\seqNiveauExposant}
% Nom court du niveau (v0.5), en texte simple et avec l'exposant.
%    \begin{macrocode}
\newcommand{\seqNiveauCourt}{6e}
\newcommand{\seqNiveauExposant}{6\textsuperscript{e}}
%    \end{macrocode}
% \end{macro}\end{macro}
%
% \subsection{Couleurs}
%    \begin{macrocode}
\definecolor{seqSixMauve}{HTML}{E8D5E8}
\definecolor{seqSixMauveBord}{HTML}{6B3D6B}
\definecolor{seqSixRouge}{HTML}{D7191C}
\definecolor{seqSixBleu}{HTML}{1F5FA8}
\definecolor{seqSixVert}{HTML}{2E8B3A}
\definecolor{seqSixPointilles}{HTML}{2E8B3A}
\definecolor{seqSixGrisTexte}{HTML}{6E6E6E}
\definecolor{seqSixEnteteTab}{HTML}{F6C9D6}
%    \end{macrocode}
%
% \subsection{Réponses}
% \begin{macro}{\ifseq@corrige}
% Le booléen qui pilote toutes les zones de réponse.
%    \begin{macrocode}
\newif\ifseq@corrige \seq@corrigefalse
\newcommand{\seqVersionCorrigee}{\seq@corrigetrue}
\newcommand{\seqVersionEleve}{\seq@corrigefalse}
\newcommand{\seqVersionComplete}{\seq@corrigetrue}
\newcommand{\seqVersionACompleter}{\seq@corrigefalse}
%    \end{macrocode}
% \end{macro}
%
% \begin{macro}{\seqTrous}\begin{macro}{\seqPointilles}
% Briques de base, sans réponse.
%    \begin{macrocode}
\NewDocumentCommand{\seqTrous}{O{3cm}}{%
  \leavevmode\hbox to #1{\leaders\hbox to 0.45em{\hss.\hss}\hfill}}
\NewDocumentCommand{\seqPointilles}{O{1} O{seqSixPointilles}}{%
  \par\begingroup\color{#2}%
  \foreach \seq@i in {1,...,#1}{%
    \noindent\leaders\hbox to 0.6em{\hss.\hss}\hfill\kern0pt\par
    \vspace{0.35\baselineskip}}%
  \endgroup}
%    \end{macrocode}
% \end{macro}\end{macro}
%
% \begin{macro}{\seqACompleter}
% En version élève, la largeur des pointillés suit celle de la réponse, pour que
% les deux versions aient la même mise en page.
%    \begin{macrocode}
\newlength{\seq@larg}
\NewDocumentCommand{\seqACompleter}{O{} m}{%
  \ifmmode\mbox{\seq@acompleter{#1}{#2}}\else\seq@acompleter{#1}{#2}\fi}
\newcommand{\seq@acompleter}[2]{%
  \ifseq@corrige{\bfseries\color{seqSixRouge}#2}%
  \else
    \if\relax\detokenize{#1}\relax
      \settowidth{\seq@larg}{\bfseries #2}%
      \seqTrous[\dimexpr\seq@larg+1.2em\relax]%
    \else\seqTrous[#1]\fi
  \fi}
%    \end{macrocode}
% \end{macro}
%
% \begin{macro}{\seqReponse}
%    \begin{macrocode}
\NewDocumentCommand{\seqReponse}{O{1} m}{%
  \ifseq@corrige
    \par{\bfseries\color{seqSixRouge}#2}\par
    \vspace{\numexpr#1-1\relax\baselineskip}%
  \else\seqPointilles[#1]\fi}
%    \end{macrocode}
% \end{macro}
%
% \begin{macro}{\seqCase}\begin{macro}{\seqCaseCocher}
%    \begin{macrocode}
\NewDocumentCommand{\seqCase}{O{1.2cm} o}{%
  \fbox{\makebox[#1]{\strut\ifseq@corrige\IfValueT{#2}{\color{seqSixRouge}#2}\fi}}}
\NewDocumentCommand{\seqCaseCocher}{o m}{%
  \ifseq@corrige\IfValueTF{#1}{{\color{seqSixRouge}$\boxtimes$}}{$\square$}%
  \else$\square$\fi\ #2\quad}
%    \end{macrocode}
% \end{macro}\end{macro}
%
% \subsection{Livret d'exercices}
% \begin{macro}{\seqCouvertureLivretSix}
% \changes{v0.6}{2026/10/19}{\cs{seqCouvertureLivretSix*} : objectifs rangés
%   par partie (\cs{seqPartieCouvertureSix}), trait de la couleur du thème}
% Forme étoilée (v0.6) : le second argument contient des
% \cs{seqPartieCouvertureSix} au lieu de la liste des objectifs.
%    \begin{macrocode}
\NewDocumentCommand{\seqCouvertureLivretSix}{s m m}{%
  \noindent\begin{minipage}[t]{0.5\linewidth}\small
    NOM : \seqTrous[4cm]\\[2pt]Prénom : \seqTrous[3.6cm]\\[2pt]
    \seqNiveauExposant\ \seqTrous[1.2cm]\end{minipage}\par
  \begin{center}{\huge\bfseries #2}\\[2pt]
    {\LARGE\bfseries Niveau \seqNiveauExposant}\end{center}
  \noindent{\bfseries\underline{Objectifs du chapitre :}}
  \IfBooleanTF{#1}{\par\smallskip#3}{%
  \begin{itemize}[itemsep=0pt,topsep=2pt]
    \foreach \seq@c/\seq@t in {#3}{\item\textit{\seq@c\ : \seq@t}}
  \end{itemize}}}
%    \end{macrocode}
% \end{macro}
%
% \begin{macro}{\seqPartieCouvertureSix}
% \changes{v0.6}{2026/10/19}{\cs{seqPartieCouvertureSix} : titre de partie et
%   objectifs de la partie sur la page de garde}
% La couleur \texttt{seqSixTheme} (bleu du module par défaut) est celle du
% thème de la séquence ; l'appli la redéfinit avant la page de garde.
%    \begin{macrocode}
\colorlet{seqSixTheme}{seqSixBleu}
\newcommand{\seqPartieCouvertureSix}[2]{%
  \par\addvspace{5pt}\noindent{\bfseries\color{seqSixTheme}#1}\par\nobreak
  \begin{tcolorbox}[enhanced, blanker, borderline west={2pt}{0pt}{seqSixTheme},
      left=8pt, top=1pt, bottom=1pt, before skip=2pt, after skip=2pt]
    \begin{itemize}[itemsep=0pt,topsep=0pt,leftmargin=1.2em]
      \foreach \seq@c/\seq@t in {#2}{\item\textit{\seq@c\ : \seq@t}}
    \end{itemize}
  \end{tcolorbox}}
%    \end{macrocode}
% \end{macro}
%
% \begin{macro}{\seqBandeauCapacite}
%    \begin{macrocode}
\newcommand{\seqBandeauCapacite}[2]{%
  \par\addvspace{0.6\baselineskip}%
  \begin{tcolorbox}[enhanced, colback=white, colframe=black, boxrule=1.2pt,
      sharp corners, halign=center, fontupper=\bfseries\itshape,
      top=3pt, bottom=3pt]
    #1 : #2\end{tcolorbox}}
%    \end{macrocode}
% \end{macro}
%
% \begin{environment}{seqActiviteSix}
%    \begin{macrocode}
\newcounter{seqActiviteSix}
\NewDocumentEnvironment{seqActiviteSix}{m}{%
  \par\addvspace{0.8\baselineskip}\stepcounter{seqActiviteSix}%
  \begin{tcolorbox}[enhanced, breakable, blanker,
      borderline west={2pt}{0pt}{seqSixBleu}, left=8pt,
      before upper={\noindent\tikz[baseline=(a.base)]\node[fill=seqSixBleu,
        text=white, font=\bfseries\small, rounded corners=2pt, inner sep=3pt](a)
        {Activité \theseqActiviteSix};\ {\bfseries #1}\par\medskip}]}{%
  \end{tcolorbox}\par\addvspace{0.5\baselineskip}}
%    \end{macrocode}
% \end{environment}
%
% \begin{environment}{seqFicheResumeSix}
%    \begin{macrocode}
\NewDocumentEnvironment{seqFicheResumeSix}{m}{%
  \begin{tcolorbox}[enhanced, breakable, colback=white, colframe=seqSixRouge,
      boxrule=1.2pt, arc=2pt, left=5pt, right=5pt,
      title={\ifseq@corrige Fiche de cours\else Fiche de résumé
             à compléter\fi\ --- #1},
      coltitle=white, colbacktitle=seqSixRouge, fonttitle=\bfseries\small]}{%
  \end{tcolorbox}}
%    \end{macrocode}
% \end{environment}
%
% \begin{macro}{\seqNumSix}\begin{environment}{seqExoSix}
%    \begin{macrocode}
\newcommand{\seqNumSix}[2]{%
  \tikz[baseline=(n.base)]\node[draw=seqSixMauveBord, fill=seqSixMauve,
    rounded corners=2pt, line width=0.8pt, minimum size=1.5em, inner sep=1pt,
    font=\bfseries](n){#1};%
  \if\relax\detokenize{#2}\relax\else\ \textbf{(#2)}\fi}
\newcounter{seqExoSix}
\NewDocumentEnvironment{seqExoSix}{O{O} o}{%
  \par\addvspace{0.8\baselineskip}\stepcounter{seqExoSix}%
  \noindent\seqNumSix{\theseqExoSix}{#1}%
  \IfValueT{#2}{\ \textit{#2}}\par\nopagebreak\vspace{2pt}%
  \begingroup\leftskip=1.9em\relax}{\par\endgroup}
%    \end{macrocode}
% \end{environment}\end{macro}
%
% \subsection{Marqueurs de l'exercice (v0.7)}
% \begin{macro}{\seqPictoOrdinateurSix}
% \changes{v0.7}{2026/10/20}{\cs{seqPictoOrdinateurSix} : écran sur son pied}
% Un écran sur son pied, à la hauteur des minuscules.
%    \begin{macrocode}
\newcommand{\seqPictoOrdinateurSix}{\tikz[baseline=0pt, x=1em, y=1em,
    line width=0.5pt, line join=round]{%
  \draw[rounded corners=0.6pt] (0,0.2) rectangle (0.95,0.8);
  \draw (0.38,0.2) -- (0.33,0.04) (0.57,0.2) -- (0.62,0.04);
  \draw (0.24,0.04) -- (0.71,0.04);}}
%    \end{macrocode}
% \end{macro}
%
% \begin{macro}{\seqPictoCalculatriceSix}
% \changes{v0.7}{2026/10/20}{\cs{seqPictoCalculatriceSix} : calculatrice, barrée
%   avec l'étoile}
% Boîtier, écran et six touches ; la forme étoilée ajoute une barre oblique épaisse.
%    \begin{macrocode}
\NewDocumentCommand{\seqPictoCalculatriceSix}{s}{\tikz[baseline=0pt, x=1em, y=1em,
    line width=0.5pt]{%
  \draw[rounded corners=0.8pt] (0,-0.05) rectangle (0.62,0.85);
  \draw (0.1,0.56) rectangle (0.52,0.75);
  \foreach \i in {0,1,2}{\foreach \j in {0,1}{%
    \fill (0.1+\i*0.16,0.06+\j*0.18) rectangle ++(0.1,0.1);}}
  \IfBooleanT{#1}{\draw[line width=1.1pt] (-0.12,-0.12) -- (0.74,0.92);}}}
%    \end{macrocode}
% \end{macro}
%
% \begin{macro}{\seqOrdinateurSix}
% \changes{v0.7}{2026/10/20}{\cs{seqOrdinateurSix} : « Ordinateur --- outil »}
%    \begin{macrocode}
\newcommand{\seqOrdinateurSix}[1]{\mbox{\seqPictoOrdinateurSix\,Ordinateur%
  \if\relax\detokenize{#1}\relax\else\ --- #1\fi}}
%    \end{macrocode}
% \end{macro}
%
% \begin{macro}{\seqCalculatriceSix}
% \changes{v0.7}{2026/10/20}{\cs{seqCalculatriceSix} : calculatrice autorisée ou
%   interdite}
%    \begin{macrocode}
\def\seq@calc@interdite{interdite}
\newcommand{\seqCalculatriceSix}[1]{\begingroup\def\seq@calc@valeur{#1}%
  \ifx\seq@calc@valeur\seq@calc@interdite
    \mbox{\seqPictoCalculatriceSix*\,calculatrice interdite}%
  \else
    \mbox{\seqPictoCalculatriceSix\,calculatrice autorisée}%
  \fi\endgroup}
%    \end{macrocode}
% \end{macro}
%
% \begin{macro}{\seqMaterielSix}
% \changes{v0.7}{2026/10/20}{\cs{seqMaterielSix} : marqueurs dans le titre d'un
%   exercice ou d'une activité}
% Remet la police droite et normale (le titre de l'exercice est en italique, celui
% de l'activité en gras).
%    \begin{macrocode}
\newcommand{\seqMaterielSix}[2]{{\normalfont\footnotesize\color{black!80}%
  \if\relax\detokenize{#1}\relax\else\seqOrdinateurSix{#1}\fi
  \if\relax\detokenize{#1}\relax\else\if\relax\detokenize{#2}\relax\else
    \quad\fi\fi
  \if\relax\detokenize{#2}\relax\else\seqCalculatriceSix{#2}\fi}}
%    \end{macrocode}
% \end{macro}
%
% \subsection{Astuce (v0.8)}
% \begin{macro}{\seqPictoAstuceSix}
% \changes{v0.8}{2026/10/20}{\cs{seqPictoAstuceSix} : ampoule au trait}
% Une ampoule au trait (globe, culot, deux bagues), à la hauteur des majuscules.
%    \begin{macrocode}
\newcommand{\seqPictoAstuceSix}{\tikz[baseline=0pt, x=1em, y=1em,
    line width=0.5pt, line join=round, line cap=round]{%
  \draw (0.36,0.22) -- (0.29,0.38) arc[start angle=235, end angle=-55,
    radius=0.28] -- (0.54,0.22) -- cycle;
  \draw (0.37,0.12) -- (0.53,0.12) (0.40,0.03) -- (0.50,0.03);
  \draw (0.40,0.30) -- (0.42,0.46) -- (0.48,0.46) -- (0.50,0.30);}}
%    \end{macrocode}
% \end{macro}
%
% \begin{macro}{\seqAstuceSix}
% \changes{v0.8}{2026/10/20}{\cs{seqAstuceSix} : astuce à la fin d'un énoncé}
% Paragraphe à part, en \cs{small}, avec un retrait suspendu de la largeur du
% pictogramme ; rien ne dépend de la version (élève ou corrigée).
%    \begin{macrocode}
\newcommand{\seqAstuceSix}[1]{\par\smallskip
  {\small\noindent\hangindent=1.3em\hangafter=1
   \makebox[1.3em][l]{\seqPictoAstuceSix}\textbf{Astuce.}\ #1\par}%
  \smallskip}
%    \end{macrocode}
% \end{macro}
%
% \subsection{Mises en route}
% \begin{environment}{seqCalculMental}
% Les items sont rangés ligne par ligne (1 2 3, puis 4 5 6), comme sur les fiches de
% l'équipe : chaque item est une minipage, et on passe à la ligne toutes les
% \meta{colonnes} minipages.
%    \begin{macrocode}
\newcounter{seqItemCM}
\newcounter{seq@cm@col}
\newcommand{\seqTitreCM}{CALCUL MENTAL}
\NewDocumentEnvironment{seqCalculMental}{m m m O{3}}{%
  \def\seq@cm@nb{#3}\def\seq@cm@cols{#4}\setcounter{seq@cm@col}{0}%
  \par\addvspace{0.5\baselineskip}\noindent
  \fbox{\parbox{\dimexpr\linewidth-2\fboxsep-2\fboxrule}%
    {\seqTitreCM\ #1 -- \textbf{#2}}}%
  \par\smallskip\noindent}{%
  \par\vspace{-0.3\baselineskip}%
  \hfill\seqScore{\seq@cm@nb}\par\addvspace{0.8\baselineskip}}
%    \end{macrocode}
% \end{environment}
% \begin{macro}{\seqItemCM}
%    \begin{macrocode}
\NewDocumentCommand{\seqItemCM}{m o}{%
  \stepcounter{seqItemCM}\stepcounter{seq@cm@col}%
  \begin{minipage}[t]{\dimexpr\linewidth/\seq@cm@cols-0.5em}
    \theseqItemCM.\ \ $#1 =$ {\color{seqSixRouge}%
      \IfValueTF{#2}{\seqACompleter[1.4cm]{$#2$}}{\seqTrous[1.4cm]}}\par
    {\color{seqSixRouge}\seqTrous[\linewidth]}\end{minipage}\hfill
  \ifnum\value{seq@cm@col}=\seq@cm@cols\relax
    \setcounter{seq@cm@col}{0}\par\medskip\noindent\fi}
%    \end{macrocode}
% \end{macro}
% \begin{macro}{\seqScore}
%    \begin{macrocode}
\newcommand{\seqScore}[1]{%
  \tikz[baseline]\node[fill=seqSixMauve!60, rounded corners=1pt, inner sep=3pt,
    font=\sffamily\bfseries\small]{SCORE : \seqTrous[1.2cm] / #1};}
%    \end{macrocode}
% \end{macro}
%
% \subsection{Planches à découper au massicot}
% La hauteur d'une bande est \cs{textheight} moins les traits de coupe, divisée par
% le nombre de bandes.
%    \begin{macrocode}
\newlength{\seq@bande}
\newcounter{seq@bande@n}
\def\seq@bande@max{1}
\newcommand{\seq@traitcoupe}{\par\noindent\makebox[\linewidth]{%
  \color{seqSixGrisTexte}\scriptsize\ding{34}\,%
  \leaders\hbox to 0.5em{\hss-\hss}\hfill}\par}
\newcommand{\seq@hauteurbande}[1]{\setlength{\seq@bande}{%
  \dimexpr(\textheight-\numexpr#1-1\relax\dimexpr2\baselineskip\relax)/#1
  -\baselineskip\relax}}
%    \end{macrocode}
% \begin{environment}{seqPlancheFiches}\begin{environment}{seqBande}
%    \begin{macrocode}
\NewDocumentEnvironment{seqPlancheFiches}{m}{%
  \clearpage\def\seq@bande@max{#1}\seq@hauteurbande{#1}%
  \setcounter{seq@bande@n}{0}}{\clearpage}
\NewDocumentEnvironment{seqBande}{}{%
  \ifnum\value{seq@bande@n}=\seq@bande@max\relax
    \clearpage\setcounter{seq@bande@n}{0}\fi
  \ifnum\value{seq@bande@n}>0 \seq@traitcoupe\fi
  \stepcounter{seq@bande@n}%
  \noindent\begin{minipage}[t][\seq@bande][t]{\linewidth}}{\end{minipage}}
%    \end{macrocode}
% \end{environment}\end{environment}
% \begin{environment}{seqCopies}
% Les compteurs du module sont mémorisés avant la première copie et rétablis
% avant chacune des suivantes.
%    \begin{macrocode}
\newcommand{\seq@sauvecpt}{\edef\seq@cpts{%
  \noexpand\setcounter{seqItemCM}{\the\value{seqItemCM}}%
  \noexpand\setcounter{seqExoSix}{\the\value{seqExoSix}}%
  \noexpand\setcounter{seqActiviteSix}{\the\value{seqActiviteSix}}%
  \noexpand\setcounter{seqExoEval}{\the\value{seqExoEval}}}}
\NewDocumentEnvironment{seqCopies}{m +b}{%
  \clearpage\seq@hauteurbande{#1}\seq@sauvecpt
  \foreach \seq@k in {1,...,#1}{\seq@cpts
    \noindent\begin{minipage}[t][\seq@bande][t]{\linewidth}#2\end{minipage}%
    \ifnum\seq@k<#1\relax\seq@traitcoupe\fi}%
  \clearpage}{}
%    \end{macrocode}
% \end{environment}
%
% \subsection{Évaluations}
%    \begin{macrocode}
\newcommand{\seqIdentEleve}[1][\seqNiveauCourt]{%
  \noindent\small Nom : \seqTrous[4cm]\hfill Prénom : \seqTrous[4cm]\hfill
  #1 \seqTrous[1cm]\hfill Date : \seqTrous[2.5cm]\par\normalsize}
\newcommand{\seqTitreEvalSix}[2]{%
  \par\medskip\begin{center}{\Large\bfseries #1}\\[2pt]
  {\small\itshape\color{seqSixGrisTexte}#2}\end{center}}
\newcommand{\seqBaremeCompetences}[2]{%
  \par\noindent\begin{tikzpicture}
    \foreach \seq@c/\seq@p [count=\seq@k] in {#1}{\xdef\seq@nb{\seq@k}}
    \pgfmathsetmacro{\seq@w}{\linewidth/1cm/(\seq@nb+1)}
    \foreach \seq@c/\seq@p [count=\seq@k from 0] in {#1,Total/#2}{%
      \draw[fill=seqSixEnteteTab] (\seq@k*\seq@w,0) rectangle ++(\seq@w,0.5)
        node[midway,font=\small\bfseries]{\seq@c};
      \draw (\seq@k*\seq@w,-0.55) rectangle ++(\seq@w,0.55)
        node[midway,font=\small]{\dots\ /\seq@p};}
  \end{tikzpicture}\par}
\newcounter{seqExoEval}
\newcommand{\seqExoEval}[2]{%
  \par\addvspace{\baselineskip}\stepcounter{seqExoEval}\noindent
  {\color{seqSixBleu}\bfseries Exercice \theseqExoEval\quad #1}\quad
  {\small\itshape\color{seqSixGrisTexte}(#2)}\par\smallskip}
\newcommand{\seqItemRond}[1]{\tikz[baseline=(c.base)]\node[draw,circle,
  inner sep=0.8pt,minimum size=1.1em,font=\small](c){#1};}
%    \end{macrocode}
%
% \subsection{Fiches de cours, travaux, activités complémentaires}
%    \begin{macrocode}
\newcommand{\seqTitreFicheSix}[1]{%
  \begin{center}{\large\bfseries\color{seqSixVert}\underline{#1}}\end{center}}
\newcommand{\seqSectionFicheSix}[1]{%
  \par\addvspace{0.4\baselineskip}\noindent
  {\bfseries\color{seqSixBleu}\underline{#1}}\par\smallskip}
\newcommand{\seqMotCle}[1]{{\bfseries\color{seqSixRouge}#1}}
\newtcolorbox{seqEncadreDefSix}{enhanced, colback=white, colframe=seqSixRouge,
  boxrule=1pt, sharp corners, left=3pt, right=3pt, top=2pt, bottom=2pt}
\NewDocumentCommand{\seqTitreTravailSix}{o m m m}{%
  \noindent\small Nom : \seqTrous[4cm]\hfill Prénom : \seqTrous[4cm]\hfill
  \seqNiveauCourt\ \seqTrous[1cm]\IfValueT{#1}{\hfill\textbf{#1}}\normalsize
  \begin{center}{\large\bfseries\color{seqSixVert}#2}\\[2pt]{\bfseries #3}\\[2pt]
  {\small\itshape #4}\end{center}}
\newtcolorbox{seqEncadreIntroSix}{enhanced, colback=white, colframe=seqSixBleu,
  boxrule=1pt, arc=2pt}
\newcommand{\seqPartieSix}[1]{\par\addvspace{0.5\baselineskip}\noindent
  {\bfseries\color{seqSixBleu}\underline{#1}}\par\smallskip}
\newcounter{seqComplSix}
\newcommand{\seqTitreComplementSix}[2]{%
  \stepcounter{seqComplSix}%
  \noindent\small Nom : \seqTrous[4.5cm]\hfill Prénom : \seqTrous[4.5cm]\hfill
  \seqNiveauCourt\ \seqTrous[1cm]\normalsize
  \begin{center}{\large\bfseries\color{seqSixVert}Activité complémentaire
    \theseqComplSix\ \ \ding{72}}\\[2pt]
  {\Large\bfseries #1}\\[2pt]
  {\small\itshape\color{seqSixGrisTexte}#2 --- pour ceux qui ont terminé}\end{center}}
%    \end{macrocode}
%
% \subsection{Figures}
% \begin{macro}{\seqDemiDroite}
%    \begin{macrocode}
\NewDocumentCommand{\seqDemiDroite}{O{12cm} m m m m O{}}{%
  \begin{tikzpicture}[x=#1/(#4)]
    \draw[thick,->] (0,0) -- (#4+0.6,0);
    \foreach \seq@i in {0,...,#4}{\draw (\seq@i,-0.12) -- (\seq@i,0.12);}
    \pgfmathtruncatemacro{\seq@sub}{#4*#5}
    \ifnum#5>1 \foreach \seq@j in {0,...,\seq@sub}{%
      \draw (\seq@j/#5,-0.06) -- (\seq@j/#5,0.06);}\fi
    \node[below] at (0,-0.12) {\small #2};
    \node[below] at (1,-0.12) {\small #3};
    \foreach \seq@x/\seq@l in {#6}{%
      \node[font=\bfseries] at (\seq@x,0.35) {\seq@l};
      \node at (\seq@x,0) {$\times$};}
  \end{tikzpicture}}
%    \end{macrocode}
% \end{macro}
% \begin{macro}{\seqFiguresSix}
% Réglages de \textsf{tkz-euclide} : points en croix, codages fins, lettres en gras.
%    \begin{macrocode}
\newcommand{\seqFiguresSix}{%
  \tkzSetUpPoint[shape=cross out, size=3pt, color=black, line width=0.6pt]%
  \tkzSetUpLine[line width=0.6pt, color=black]%
  \tkzSetUpLabel[font=\bfseries\small]}
\AtBeginDocument{\@ifpackageloaded{tkz-euclide}{\seqFiguresSix}{}}
%    \end{macrocode}
% \end{macro}
%
% \subsection{Macros issues des premiers lots (v0.4)}
% Ces macros ont été écrites pendant la conversion des séquences S01 à S03 et
% gardent leur préfixe \cs{seqLoc} pour que les énoncés déjà convertis compilent
% sans retouche. Elles suivent toutes la bascule élève / corrigé :
% \begin{itemize}
% \item \cs{seqLocSiCorrige}\marg{corrigé}\marg{élève}, \cs{seqLocCellule}\marg{réponse},
%   \cs{seqLocEntoure}\oarg{couleur}\marg{texte}, \cs{seqLocExplication}\marg{texte} ;
% \item tableaux de numération colorés : \cs{seqLocTableauNumeration}, \cs{seqLocTableauDecimal} ;
% \item hiéroglyphes : \cs{seqLocHiero}, \cs{seqLocHieroNombre} ; bande de fractions : \cs{seqLocBande} ;
% \item demi-droites : \cs{seqLocCroix}, \cs{seqLocSurligne} ;
% \item opérations posées : \cs{seqLocChiffres}, \cs{seqLocTrait}, \cs{seqLocGrilleOp},
%   \cs{seqLocPotence}, \cs{seqLocPotenceVisible} ;
% \item ordre de grandeur : \cs{seqLocFlechesOG}, \cs{seqLocTableauOG} ; calculatrice : \cs{seqLocEcran}.
% \end{itemize}
%
%    \begin{macrocode}
\usetikzlibrary{shapes.geometric}

\usetikzlibrary{matrix}



%% --- Bascule élève / corrigé (enveloppe publique de \ifseq@corrige) ---------
%% \seqLocSiCorrige{version corrigée}{version élève}
\newcommand{\seqLocSiCorrige}[2]{\ifseq@corrige#1\else#2\fi}

%% --- Réponse sans pointillés, pour une cellule de tableau ou une case -------
%% \seqLocCellule{réponse} : vide (même largeur) en version élève.
\newcommand{\seqLocCellule}[1]{\leavevmode\seqLocSiCorrige
  {{\bfseries\color{seqSixRouge}#1}}{\phantom{\bfseries #1}}}

%% --- « Entoure » : la réponse est entourée en rouge dans le corrigé ---------
%% \seqLocEntoure{texte} : à utiliser seulement pour les bonnes réponses.
%% \seqLocEntoure[couleur]{texte} (couleur par défaut : seqSixRouge).
\NewDocumentCommand{\seqLocEntoure}{O{seqSixRouge} m}{\ifseq@corrige\def\seq@entc{#1}\else
  \def\seq@entc{none}\fi\tikz[baseline=(e.base)]\node[ellipse,
  inner sep=1.5pt, draw=\seq@entc, line width=0.9pt](e){#2};}

%% --- Explications du corrigé (champ « corrige » de l'appli) ---------------
\newcommand{\seqLocExplication}[1]{\ifseq@corrige\par\smallskip
  \begin{tcolorbox}[enhanced, breakable, width=\dimexpr\linewidth-\leftskip\relax, colback=seqSixMauve!25,
    colframe=seqSixMauveBord!60, boxrule=0.5pt, arc=1pt, left=4pt, right=4pt,
    top=2pt, bottom=2pt, fontupper=\small]
  \textbf{Explications.} #1\end{tcolorbox}\fi}

%% --- Tableau de numération coloré (gabarit des livrets 2026 de l'équipe) ---
\definecolor{seqLocMillions}{HTML}{F7A8F7}
\definecolor{seqLocMilliers}{HTML}{5BEB86}
\definecolor{seqLocUnites}{HTML}{72CDFA}
\definecolor{seqLocDecimale}{HTML}{FFF0A0}
\newcommand{\seqLoc@rot}[1]{\rotatebox{90}{\small\bfseries #1\ }}
%% \seqLocTableauNumeration{lignes} : 12 colonnes (milliards -> unités).
%% Chaque ligne : 12 cellules séparées par &, terminée par \\ \hline.
\newcommand{\seqLocTableauNumeration}[1]{%
  \begin{center}\renewcommand{\arraystretch}{1.25}\small
  \begin{tabular}{|*{3}{>{\centering\arraybackslash}p{0.62cm}|}*{3}{>{\columncolor{seqLocMillions!45}\centering\arraybackslash}p{0.62cm}|}%
      *{3}{>{\columncolor{seqLocMilliers!45}\centering\arraybackslash}p{0.62cm}|}*{3}{>{\columncolor{seqLocUnites!45}\centering\arraybackslash}p{0.62cm}|}}
  \hline
  \multicolumn{3}{|c|}{\bfseries Milliards}&
  \multicolumn{3}{c|}{\cellcolor{seqLocMillions}\bfseries Millions}&
  \multicolumn{3}{c|}{\cellcolor{seqLocMilliers}\bfseries Milliers}&
  \multicolumn{3}{c|}{\cellcolor{seqLocUnites}\bfseries Unités}\\ \hline
  \seqLoc@rot{centaines}&\seqLoc@rot{dizaines}&\seqLoc@rot{unités}&
  \seqLoc@rot{centaines}&\seqLoc@rot{dizaines}&\seqLoc@rot{unités}&
  \seqLoc@rot{centaines}&\seqLoc@rot{dizaines}&\seqLoc@rot{unités}&
  \seqLoc@rot{centaines}&\seqLoc@rot{dizaines}&\seqLoc@rot{unités}\\ \hline
  #1
  \end{tabular}\end{center}}
%% \seqLocTableauDecimal{lignes} : classe des milliers, classe des unités,
%% virgule, dixièmes, centièmes, millièmes (10 cellules par ligne).
\newcommand{\seqLocTableauDecimal}[1]{%
  \begin{center}\renewcommand{\arraystretch}{1.25}\small
  \begin{tabular}{|*{3}{>{\columncolor{seqLocMilliers!45}\centering\arraybackslash}p{0.62cm}|}%
      *{3}{>{\columncolor{seqLocUnites!45}\centering\arraybackslash}p{0.62cm}|}>{\centering\arraybackslash}p{0.2cm}|%
      *{3}{>{\columncolor{seqLocDecimale!60}\centering\arraybackslash}p{0.62cm}|}}
  \hline
  \multicolumn{3}{|c|}{\cellcolor{seqLocMilliers}\bfseries Milliers}&
  \multicolumn{3}{c|}{\cellcolor{seqLocUnites}\bfseries Unités}& &
  \multicolumn{3}{c|}{\cellcolor{seqLocDecimale}\bfseries Partie décimale}\\ \hline
  \seqLoc@rot{centaines}&\seqLoc@rot{dizaines}&\seqLoc@rot{unités}&
  \seqLoc@rot{centaines}&\seqLoc@rot{dizaines}&\seqLoc@rot{unités}& &
  \seqLoc@rot{dixièmes}&\seqLoc@rot{centièmes}&\seqLoc@rot{millièmes}\\ \hline
  #1
  \end{tabular}\end{center}}

%% --- Hiéroglyphes égyptiens (images extraites du TP1 de l'équipe) ----------
%% \seqLocHiero{k}{n} : n fois le symbole de valeur 10^k (k = 0 à 6).
\newcommand{\seqLocHiero}[2]{\ifnum#2>0\relax\foreach \seq@h in {1,...,#2}{%
  \includegraphics[height=1.15em]{hiero-#1.png}\hspace{1pt}}\fi}
%% \seqLocHieroNombre{millions}{c. de mille}{d. de mille}{milliers}{centaines}{dizaines}{unités}
\newcommand{\seqLocHieroNombre}[7]{\mbox{\raisebox{-0.25em}{%
  \seqLocHiero{6}{#1}\seqLocHiero{5}{#2}\seqLocHiero{4}{#3}\seqLocHiero{3}{#4}%
  \seqLocHiero{2}{#5}\seqLocHiero{1}{#6}\seqLocHiero{0}{#7}}}}

%% --- Bande unité partagée (fractions) ---------------------------------------
%% \seqLocBande[longueur]{nombre de parts}{parts coloriées}
\NewDocumentCommand{\seqLocBande}{O{6cm} m m}{%
  \begin{tikzpicture}[x=#1/#2]
    \ifnum#3>0\relax\fill[seqSixBleu!25] (0,0) rectangle (#3,0.5);\fi
    \foreach \seq@b in {1,...,#2}{\draw (\seq@b-1,0) rectangle (\seq@b,0.5);}
    \draw[thick] (0,0) rectangle (#2,0.5);
  \end{tikzpicture}}

%% --- Demi-droites graduées (S02) -------------------------------------------
%% \seqLocCroix{x}{nom} : point marqué d'une croix et nommé au-dessus (dans un tikzpicture).
\newcommand{\seqLocCroix}[2]{\node at (#1,0) {$\times$};
  \node[above,font=\bfseries] at (#1,0.14) {#2};}
%% \seqLocSurligne{x1}{x2} : surligne la partie [x1 ; x2] d'un axe (zooms successifs).
\newcommand{\seqLocSurligne}[2]{\fill[seqSixBleu!45] (#1,-0.12) rectangle (#2,0.12);}

%% --- Opérations posées (S03) -------------------------------------------------
%% Repère commun : une colonne de chiffres = \seqLocCol, une ligne = \seqLocLig ;
%% la ligne k est à l'ordonnée k (vers le bas), la colonne c occupe [c ; c+1].
\newlength{\seqLocCol}\setlength{\seqLocCol}{0.34cm}
\newlength{\seqLocLig}\setlength{\seqLocLig}{0.5cm}
\newcount\seqLoc@x
%% \seqLocChiffres{colonne}{ligne}{chaîne} (dans un tikzpicture de ce repère) :
%% un caractère par colonne, à partir de la colonne donnée. Codes : _ case vide,
%% * signe fois, - signe moins, + signe plus ; la virgule ne prend pas de colonne : elle est
%% placée en bas à droite du chiffre qui la précède.
\def\seqLoc@fin{\seqLoc@fin}
\newcommand{\seqLocChiffres}[3]{\def\seqLoc@y{#2}\seqLoc@x=#1\relax
  \seqLoc@parse#3\seqLoc@fin}
\def\seqLoc@parse#1{\ifx#1\seqLoc@fin\else\seqLoc@car{#1}\expandafter\seqLoc@parse\fi}
\def\seqLoc@car#1{%
  \if#1,\node[inner sep=0pt,anchor=north] at (\the\seqLoc@x,\seqLoc@y+0.12) {,};%
  \else\if#1_\advance\seqLoc@x by 1\relax
  \else\node[inner sep=0pt] at (\the\seqLoc@x+0.5,\seqLoc@y) {\strut\seqLoc@sym{#1}};%
    \advance\seqLoc@x by 1\relax\fi\fi}
\def\seqLoc@sym#1{\if#1*$\times$\else\if#1-$-$\else\if#1+$+$\else#1\fi\fi\fi}
%% \seqLocTrait{c1}{c2}{ligne} : trait d'opération sous la ligne, de la colonne c1 à c2 (incluses).
\newcommand{\seqLocTrait}[3]{\draw[line width=0.9pt] (#1,#3+0.5) -- (#2+1,#3+0.5);}
%% \seqLocGrilleOp{colonnes}{lignes}{visible}{corrigé} : opération posée sur un
%% quadrillage (gabarit des grilles de multiplication de l'équipe). Le contenu
%% visible (facteurs) est toujours affiché ; le corrigé (en rouge) seulement
%% dans la version corrigée.
\NewDocumentCommand{\seqLocGrilleOp}{m m m m}{%
  \begin{tikzpicture}[x=\seqLocCol,y=-\seqLocLig,baseline={(0,0)}]
    \begin{scope}[seqSixBleu!35,line width=0.4pt]
      \foreach \seqLoc@i in {0,...,#1}{\draw (\seqLoc@i,-0.5) -- (\seqLoc@i,#2-0.5);}
      \foreach \seqLoc@j in {0,...,#2}{\draw (0,\seqLoc@j-0.5) -- (#1,\seqLoc@j-0.5);}
    \end{scope}
    #3
    \seqLocSiCorrige{\begin{scope}[color=seqSixRouge,font=\bfseries]#4\end{scope}}{}
  \end{tikzpicture}}
%% \seqLocPotence[lignes]{colonnes}{dividende}{diviseur}{quotient}{corrigé}
%% Potence de division. Le dividende (chaîne de \seqLocChiffres) commence à la
%% colonne 1 ; la barre verticale est après la colonne <colonnes>. Le corrigé
%% (étapes, en rouge) et le quotient n'apparaissent que dans la version corrigée ;
%% en version élève, la place est réservée (<lignes> lignes, 5 par défaut).
\newlength{\seqLoc@la}\newlength{\seqLoc@lb}
\NewDocumentCommand{\seqLocPotence}{O{5} m m m m m}{%
  \settowidth{\seqLoc@la}{#4}\settowidth{\seqLoc@lb}{\bfseries #5}%
  \ifdim\seqLoc@lb>\seqLoc@la\seqLoc@la=\seqLoc@lb\fi
  \addtolength{\seqLoc@la}{1.2em}%
  \begin{tikzpicture}[x=\seqLocCol,y=-\seqLocLig,baseline={(0,0.1)}]
    \seqLocChiffres{1}{0}{#3}
    \draw[line width=0.9pt] (#2+0.4,-0.5) -- (#2+0.4,#1+0.4);
    \draw[line width=0.9pt] (#2+0.4,0.5) -- ++(\the\seqLoc@la,0);
    \node[anchor=west,inner sep=0pt] at (#2+0.8,0) {\strut #4};
    \ifseqLoc@visible
      \node[anchor=west,inner sep=0pt] at (#2+0.8,1) {\strut #5};#6
    \else
    \seqLocSiCorrige{\node[anchor=west,inner sep=0pt,text=seqSixRouge,font=\bfseries]
      at (#2+0.8,1) {\strut #5};
      \begin{scope}[color=seqSixRouge,font=\bfseries]#6\end{scope}}{}\fi
  \end{tikzpicture}}
%% \seqLocPotenceVisible : mêmes arguments, mais tout est affiché (en noir) dans les
%% deux versions : division donnée en exemple, ou division à analyser.
\newif\ifseqLoc@visible
\NewDocumentCommand{\seqLocPotenceVisible}{O{5} m m m m m}{{\seqLoc@visibletrue
  \seqLocPotence[#1]{#2}{#3}{#4}{#5}{#6}}}

%% --- Ordre de grandeur (S03) -------------------------------------------------
%% \seqLocFlechesOG{terme 1}{opération}{terme 2}{arrondi 1}{arrondi 2}{résultat}
%% Calcul en haut, flèches vers les nombres « ronds » choisis, ordre de grandeur du
%% résultat. Les trois réponses passent par \seqACompleter.
\NewDocumentCommand{\seqLocFlechesOG}{m m m m m m}{%
  \begin{tikzpicture}[baseline={(0,0)},every node/.style={inner sep=1pt}]
    \node (a) at (0,0) {\strut #1}; \node at (1.35,0) {\strut #2};
    \node (b) at (2.7,0) {\strut #3}; \node at (3.9,0) {\strut $=$};
    \node (c) at (4.9,0) {\strut ?};
    \node (a2) at (0,-1) {\strut\seqACompleter[1.4cm]{#4}};
    \node at (1.35,-1) {\strut #2};
    \node (b2) at (2.7,-1) {\strut\seqACompleter[1.4cm]{#5}};
    \node at (3.9,-1) {\strut $\approx$};
    \node (c2) at (4.9,-1) {\strut\seqACompleter[1.4cm]{#6}};
    \foreach \seqLoc@h/\seqLoc@b in {a/a2,b/b2,c/c2}{\draw[->,seqSixBleu,thick]
      (\seqLoc@h.south) -- (\seqLoc@b.north);}
    \draw[seqSixGrisTexte,rounded corners=3pt] (-1.2,0.45) rectangle (5.6,-1.45);
  \end{tikzpicture}}
%% \seqLocTableauOG{lignes} : tableau « Opération / Ordre de grandeur des termes /
%% Ordre de grandeur du résultat / Résultat exact ». Cellules séparées par \&,
%% lignes terminées par \\ ; les réponses passent par \seqLocCellule.
\newcommand{\seqLocTableauOG}[1]{%
  \begin{tikzpicture}
    \matrix[matrix of nodes, ampersand replacement=\&, nodes in empty cells,
      column sep=-\pgflinewidth, row sep=-\pgflinewidth,
      nodes={draw, minimum height=0.85cm, align=center, anchor=center, inner sep=2pt, font=\small},
      column 1/.style={nodes={text width=3.3cm}},
      column 2/.style={nodes={text width=3.6cm}},
      column 3/.style={nodes={text width=2.6cm}},
      column 4/.style={nodes={text width=2.6cm}},
      row 1/.style={nodes={fill=seqSixEnteteTab, font=\small\bfseries, minimum height=1.25cm}}]{
      Opération \& Ordre de grandeur des termes \& Ordre de grandeur du résultat \& Résultat exact\\
      #1};
  \end{tikzpicture}}

%% --- Écran de calculatrice (S03) ---------------------------------------------
%% \seqLocEcran{calcul} : petit écran affichant un calcul (chiffres cachés, etc.).
\newcommand{\seqLocEcran}[1]{\tikz[baseline=(e.base)]\node[draw=seqSixGrisTexte,
  line width=0.8pt, rounded corners=3pt, fill=seqSixBleu!8, inner sep=5pt,
  font=\ttfamily\large](e){#1};}
%    \end{macrocode}
%
% \subsection{Géométrie plane (v0.6)}
% Macros écrites pendant la conversion de la séquence S05 ; leurs noms
% \cs{seqLoc} sont ceux des énoncés du lot.
%
% \begin{macro}{\seqLocCadreTrace}
% \changes{v0.6}{2026/10/19}{\cs{seqLocCadreTrace} : zone de tracé, vide en
%   version élève, avec la figure en version corrigée (lot S05)}
% Le cadre a la même hauteur dans les deux versions ; \cs{mbox} garde la boîte
% non vide en version élève.
%    \begin{macrocode}
\NewDocumentCommand{\seqLocCadreTrace}{O{5cm} m}{\par\smallskip
  \begin{tcolorbox}[enhanced, width=\dimexpr\linewidth-\leftskip\relax,
      height=#1, colback=white, colframe=seqSixGrisTexte!40, boxrule=0.4pt,
      arc=1pt, valign=center, halign=center, left=2pt, right=2pt, top=2pt,
      bottom=2pt]%
    \ifseq@corrige #2\else\mbox{}\fi
  \end{tcolorbox}}
%    \end{macrocode}
% \end{macro}
%
% \begin{macro}{\seqLocAngle}
% \changes{v0.6}{2026/10/19}{\cs{seqLocAngle} : angle géométrique avec
%   chapeau (lot S05)}
%    \begin{macrocode}
\newcommand{\seqLocAngle}[1]{\ensuremath{\widehat{#1}}}
%    \end{macrocode}
% \end{macro}
%
% Style TikZ \texttt{seqLocCorr} des tracés du corrigé.
% \changes{v0.6}{2026/10/19}{Style TikZ \texttt{seqLocCorr} des tracés du
%   corrigé (lot S05)}
%    \begin{macrocode}
\tikzset{seqLocCorr/.style={color=seqSixRouge, line width=0.8pt}}
%    \end{macrocode}
%
% \begin{macro}{\seqLocRapporteur}
% \changes{v0.6}{2026/10/19}{\cs{seqLocRapporteur} : demi-rapporteur
%   semi-transparent à double graduation, à la place de \cs{tkzProtractor}
%   (lot S05)}
% Le demi-disque est tracé dans un \texttt{scope} déplacé au centre et tourné
% de l'orientation : graduations de 1 en 1, de 5 en 5 et de 10 en 10, nombres
% de l'échelle extérieure (0 à 180) puis de l'échelle intérieure (180 à 0).
%    \begin{macrocode}
\NewDocumentCommand{\seqLocRapporteur}{O{3} m m}{%
  \begin{scope}[shift={#2}, rotate=#3]
    \fill[seqSixBleu!8, fill opacity=0.75] (#1,0) arc (0:180:#1) -- cycle;
    \draw[seqSixGrisTexte, line width=0.4pt] (-#1-0.15,0) -- (#1+0.15,0) (#1,0) arc (0:180:#1);
    \draw[seqSixGrisTexte, line width=0.3pt] (#1*0.55,0) arc (0:180:#1*0.55);
    \draw[seqSixGrisTexte, line width=0.25pt] \foreach \seq@a in {0,...,180}{(\seq@a:#1) -- (\seq@a:#1-0.08)}
      \foreach \seq@a in {0,5,...,180}{(\seq@a:#1-0.08) -- (\seq@a:#1-0.16)}
      \foreach \seq@a in {0,10,...,180}{(\seq@a:#1-0.16) -- (\seq@a:#1-0.24)};
    \foreach \seq@a in {0,10,...,180}{%
      \pgfmathtruncatemacro{\seq@b}{180-\seq@a}%
      \node[font=\fontsize{4.5}{5}\selectfont, rotate=\seq@a-90, inner sep=0pt] at (\seq@a:#1-0.36) {\seq@a};
      \node[font=\fontsize{4.5}{5}\selectfont, rotate=\seq@a-90, inner sep=0pt, seqSixRouge!80!black] at (\seq@a:#1-0.62) {\seq@b};}
    \draw[seqSixGrisTexte, line width=0.4pt] (0,0) -- (0,0.22) (-0.12,0) -- (0.12,0);
  \end{scope}}
%    \end{macrocode}
% \end{macro}
%
% \subsection{Pré-algèbre (v0.6)}
% Macros écrites pendant la conversion de la séquence S04 ; leurs noms
% \cs{seqLoc} sont ceux des énoncés du lot.
%
% Styles TikZ des schémas en barres.
% \changes{v0.6}{2026/10/19}{Styles TikZ \texttt{seqLocBarre}, \texttt{seqLocBarreR},
%   \texttt{seqLocBarreV}, \texttt{seqLocBarreQ} des schémas en barres (lot S04)}
%    \begin{macrocode}
\usetikzlibrary{decorations.pathreplacing}
\tikzset{seqLocBarre/.style={draw=seqSixBleu!80!black, fill=seqSixBleu!12, line width=0.6pt},
  seqLocBarreR/.style={seqLocBarre, draw=seqSixRouge!80!black, fill=seqSixRouge!10},
  seqLocBarreV/.style={seqLocBarre, draw=seqSixVert!80!black, fill=seqSixVert!14},
  seqLocBarreQ/.style={seqLocBarre, fill=white, dashed}}
%    \end{macrocode}
%
% \begin{macro}{\seqLocBloc}
% \changes{v0.6}{2026/10/19}{\cs{seqLocBloc} : bloc de schéma en barres (lot S04)}
%    \begin{macrocode}
\NewDocumentCommand{\seqLocBloc}{O{seqLocBarre} m m m m}{%
  \draw[#1] (#2,#3) rectangle ++(#4,0.6);
  \node[font=\small, inner sep=1pt] at ($(#2,#3)+(#4/2,0.3)$) {#5};}
%    \end{macrocode}
% \end{macro}
%
% \begin{macro}{\seqLocAccolade}
% \changes{v0.6}{2026/10/19}{\cs{seqLocAccolade} : accolade horizontale sous ou
%   sur une barre (lot S04)}
%    \begin{macrocode}
\NewDocumentCommand{\seqLocAccolade}{O{below} m m m m}{%
  \def\seq@acc@pos{#1}\def\seq@acc@below{below}%
  \ifx\seq@acc@pos\seq@acc@below
    \draw[decorate, decoration={brace, amplitude=5pt, mirror}, line width=0.6pt]
      (#2,#4) -- (#3,#4) node[midway, below=5pt, font=\small, inner sep=1pt] {#5};
  \else
    \draw[decorate, decoration={brace, amplitude=5pt}, line width=0.6pt]
      (#2,#4) -- (#3,#4) node[midway, above=5pt, font=\small, inner sep=1pt] {#5};
  \fi}
%    \end{macrocode}
% \end{macro}
%
% \begin{macro}{\seqLocBalance}
% \changes{v0.6}{2026/10/19}{\cs{seqLocBalance} : balance de Roberval en
%   équilibre (lot S04)}
%    \begin{macrocode}
\NewDocumentCommand{\seqLocBalance}{O{2.4} m m}{%
  \begin{tikzpicture}[baseline=(seqbal.base)]
    \fill[seqSixGrisTexte!45] (-0.7,0) rectangle (0.7,0.15);
    \fill[seqSixGrisTexte!70] (-0.18,0.15) -- (0.18,0.15) -- (0.06,1.05) -- (-0.06,1.05) -- cycle;
    \draw[line width=1.8pt, seqSixGrisTexte] (-#1,1.05) -- (#1,1.05);
    \fill[seqSixGrisTexte] (0,1.05) circle (2.2pt);
    \draw[line width=1.2pt, seqSixGrisTexte] (-#1,1.05) -- (-#1,1.35) (#1,1.05) -- (#1,1.35);
    \draw[line width=2.2pt, seqSixGrisTexte] (-#1-0.45*#1-0.3,1.38) -- (-#1+0.45*#1+0.3,1.38)
      (#1-0.45*#1-0.3,1.38) -- (#1+0.45*#1+0.3,1.38);
    \node[anchor=south, inner sep=1pt] at (-#1,1.45) {#2};
    \node[anchor=south, inner sep=1pt] at (#1,1.45) {#3};
    \node[inner sep=0pt] (seqbal) at (0,0) {\strut};
  \end{tikzpicture}}
%    \end{macrocode}
% \end{macro}
%
% \begin{macro}{\seqLocSac}
% \changes{v0.6}{2026/10/19}{\cs{seqLocSac} : sac de masse inconnue (lot S04)}
%    \begin{macrocode}
\newcommand{\seqLocSac}[1]{\tikz[baseline=(s.base)]\node[draw=seqSixBleu!80!black,
  fill=seqSixBleu!18, rounded corners=3pt, minimum width=0.62cm, minimum height=0.62cm,
  inner sep=1pt, font=\small\bfseries](s){#1};}
%    \end{macrocode}
% \end{macro}
%
% \begin{macro}{\seqLocMasse}
% \changes{v0.6}{2026/10/19}{\cs{seqLocMasse} : masse marquée (lot S04)}
%    \begin{macrocode}
\newcommand{\seqLocMasse}[1]{\tikz[baseline=(m.base)]\node[draw=black, fill=seqSixGrisTexte!25,
  trapezium, trapezium angle=70, minimum height=0.5cm, inner sep=1.5pt,
  font=\scriptsize\bfseries](m){#1};}
%    \end{macrocode}
% \end{macro}
%
% \begin{macro}{\seqLocObjet}
% \changes{v0.6}{2026/10/19}{\cs{seqLocObjet} : objet rond posé sur la balance
%   (lot S04)}
%    \begin{macrocode}
\NewDocumentCommand{\seqLocObjet}{O{seqSixRouge} m}{\tikz[baseline=(o.base)]\node[draw=#1!80!black,
  fill=#1!25, circle, minimum size=0.62cm, inner sep=0.5pt, font=\scriptsize\bfseries](o){#2};}
%    \end{macrocode}
% \end{macro}
%
%    \begin{macrocode}
%</package>
%    \end{macrocode}
% \Finale
\endinput
